% Transcription — fonds Alexandre Grothendieck, shelfmark 159, batch 1 (pages 1-20).
% Established from the facsimile archives/batches/159/batch-01.pdf (a mirror of
% https://grothendieck.umontpellier.fr/159.pdf), per the skill
% transcribe-grothendieck. No cover sheet: PDF page N = archivists' page N,
% checked against the pencilled numbers bottom-left.
% Page 13 carries nothing but the archivists' number — absent.
% Pass: Opus 5 (claude-opus-5), 2026-09-19 — first pass, unchecked against the pages by a human.
\documentclass[11pt,a4paper]{article}
% Le préambule commun à toutes les transcriptions du fonds Grothendieck.
%
% Il tient en une idée : **l'incertitude se compose, elle ne se résout pas.**
% Une transcription qui lisse un mot douteux a détruit la seule chose pour
% laquelle on la faisait — un lecteur doit pouvoir savoir, à chaque endroit,
% ce qui était sur le feuillet et ce qui a été ajouté. Les six macros qui
% suivent sont donc l'appareil critique, et non de la décoration.
%
% Les mêmes commandes sont reconnues par `scripts/render.mjs`, qui produit la
% vue de lecture du site. Une macro ajoutée ici doit l'être aussi là-bas,
% faute de quoi le rendu échouera bruyamment — ce qui est voulu : un rendu
% silencieusement faux est pire qu'un rendu absent.

\ProvidesPackage{grothendieck}[2026/08/08 Transcriptions du fonds Grothendieck]

\usepackage{amsmath,amssymb,amsthm}
\usepackage{mathrsfs}
\usepackage{stmaryrd}
% Les diagrammes commutatifs sont partout dans ce fonds : c'est la langue
% même de la géométrie algébrique de Grothendieck, et une transcription qui
% les rendrait en prose ne transcrirait rien.
\usepackage{tikz-cd}
% Français par défaut — c'est la langue des feuillets — mais l'anglais est
% chargé aussi : la traduction et le résumé sont des documents anglais, et la
% typographie française (espaces avant les deux-points) y serait une faute.
% Ces documents basculent par \selectlanguage{english} en tête de corps.
\usepackage[english,french]{babel}
\usepackage{fontspec}
\usepackage{geometry}
\geometry{a4paper,margin=25mm,marginparwidth=28mm}
\usepackage{marginnote}
\usepackage{xcolor}
% ulem, not soul. `soul`'s \st and \ul are built on a character-by-character
% scan that XeLaTeX's font machinery breaks: the document fails with an
% undefined control sequence rather than degrading. ulem does the same two jobs
% and is Unicode-engine safe. `normalem` stops it hijacking \emph, which the
% transcriptions use for Grothendieck's own emphasis.
\usepackage[normalem]{ulem}
\usepackage{hyperref}
\hypersetup{colorlinks=true,linkcolor=black,urlcolor=black,pdfborder={0 0 0}}

% --- Métadonnées du lot -----------------------------------------------------
% Reprises telles quelles par le script de rendu, qui les affiche en tête de
% la vue de lecture. Elles fixent ce que le fichier prétend couvrir : sans
% elles, une transcription flotte sans point d'attache dans un fonds de seize
% mille feuillets.

\newcommand{\folder}[1]{\gdef\grothFolder{#1}}
\newcommand{\batch}[1]{\gdef\grothBatch{#1}}
\newcommand{\leaves}[2]{\gdef\grothFirst{#1}\gdef\grothLast{#2}}
\newcommand{\foldertitle}[1]{\gdef\grothFoldertitle{#1}}
\gdef\grothFolder{?}\gdef\grothBatch{?}\gdef\grothFirst{?}\gdef\grothLast{?}\gdef\grothFoldertitle{}

% --- Le résumé --------------------------------------------------------------
%
% Il ouvre la lecture modernisée, sous le titre « Résumé », et il est composé
% en petit. Ce n'est pas un ornement : c'est le seul passage du document qui
% ne soit pas des mathématiques, et un lecteur doit pouvoir le reconnaître
% avant de l'avoir lu — puis le sauter s'il connaît déjà le sujet, ou s'y
% arrêter s'il ne le connaît pas. Le filet qui le suit marque la reprise du
% corps.

\newenvironment{resume}
  {\small\setlength{\parskip}{0.45em}}
  {\par\medskip\hrule\medskip}

% --- Mots-clés ---------------------------------------------------------------
%
% En anglais, en clôture du résumé de la lecture modernisée : le vocabulaire
% moderne sous lequel chercher ce que les feuillets construisent. Le manifeste
% du site (`npm run manifest`) les extrait pour étiqueter la cote entière —
% cette ligne est la source unique des tags, il n'y a pas de fichier à côté.
\newcommand{\keywords}[1]{\par\smallskip\noindent
  {\footnotesize\textsc{Keywords} --- \textit{#1}}\par}

% --- L'appareil critique ----------------------------------------------------

% Le numéro de feuillet, en marge. C'est l'ancre qui permet de régler un
% désaccord sur une formule en regardant *un* feuillet plutôt qu'une chemise
% de six cents. Le site s'en sert aussi pour tourner les pages du fac-similé
% quand on fait défiler la transcription.
\newcommand{\leaf}[1]{%
  \marginnote{\footnotesize\color{gray!70}\textsf{#1}}%
  \label{leaf:#1}\ignorespaces}

% Illisible : on ne devine pas. Un mot inventé qui se lit comme les autres est
% le pire résultat possible.
\newcommand{\ill}{\textcolor{red!60!black}{[\,\ldots\,]}}

% Lecture douteuse : le mot est donné, et signalé comme incertain.
\newcommand{\uncertain}[1]{\uline{#1}}

% Ajout de l'éditeur — une lettre manquante, un mot sous-entendu.
\newcommand{\add}[1]{\textcolor{blue!50!black}{[#1]}}

% Biffé par Grothendieck lui-même. Ce qu'il a rayé fait partie du document :
% une rature dit souvent où la pensée a changé de direction.
\newcommand{\struck}[1]{\sout{#1}}

% Note du transcripteur, sur l'état matériel du feuillet.
\newcommand{\note}[1]{\par\smallskip{\small\color{gray!60!black}\textit{[#1]}}\par\smallskip}

% Note marginale de Grothendieck, distincte du corps du texte.
\newcommand{\marginal}[1]{\par\smallskip\noindent\hspace{1em}%
  {\small\color{gray!60!black}#1}\par\smallskip}

% --- Titre ------------------------------------------------------------------

\AtBeginDocument{%
  \begin{center}
    {\large\bfseries Fonds Alexandre Grothendieck --- cote n\textsuperscript{o}~\grothFolder}\\[2pt]
    {\itshape\grothFoldertitle}\\[4pt]
    {\small feuillets \grothFirst--\grothLast\ (lot \grothBatch)}\\[2pt]
    {\footnotesize\color{gray!60!black}Transcription établie d'après le fac-similé numérisé
     par l'Université de Montpellier.\\
     Les lectures incertaines sont soulignées, les illisibles notées [\,\ldots\,],
     les ajouts entre crochets.}
  \end{center}
  \vspace{1em}}

\endinput


\folder{159}
\batch{1}
\pages{1}{20}
\dating{s.d. — le groupe « Dérivateurs » (157-1 à 160) est daté 1990-[1991]}
\watermark{Édition de démonstration}
\foldertitle{Connexions projectives et uniformisation : notes manuscrites (s.d.)}

\begin{document}

\section*{Connexions projectives canoniques et uniformisation des courbes algébriques}

\page{1}
\note{le titre ci-dessus est de lui, porté en tête du feuillet}

\subsection*{1. Connexions projectives}

a) $S$ schéma de base (disons). La donnée d'un fibré de Brauer-Séveri $P$ sur
$S$ (schéma propre et lisse $/S$ dont les fibres sont des espaces projectifs) de
dim. relative constante $r$ revient à : la donnée d'un torseur $T$ sur $S$ de
groupe $\mathrm{PGL}(r)_S$. La donnée d'un $P$ avec section \ill{} revient à se
donner une restriction du groupe structural de $\mathrm{PGL}(r)_S$ au
sous-groupe qui fixe une section donnée de \struck{$P$}
$P^r_0 = P(\underline{O}_S^{\,r+1})$, disons celle qui correspond au quotient
$\underline{O}_S$ de $\underline{O}_S^{\,r+1}$ « correspondant à la dernière
coordonnée ». Ce sous-groupe s'identifie au groupe affine $\mathrm{Aff}(r)_S$,
produit semi-direct
\[
  \mathrm{Aff}(r) \;=\; \mathrm{GL}(r)\cdot E^{r}
\]
\note{dans la formule, l'argument de $\mathrm{Aff}$ et celui de $\mathrm{GL}$
sont chacun écrits sur une lettre biffée}
où $\mathrm{GL}(r)$ agit sur $E^r$ de la façon évidente.

[ \struck{Un} Groupe des \struck{auto}morphismes du module
$\underline{O}_S^{\,r+1} = \underline{O}_S^{\,r} \times \underline{O}_S
= E_0 \times L_0$ qui respecte la structure d'extension, \struck{\ill{}} est en
effet celui des matrices
\[
  {\scriptstyle r+1}
  \left\{
  \begin{pmatrix} A & C \\ 0 & 1 \end{pmatrix}
  \right\}
  {\scriptstyle r}
\]
avec $A$ \struck{\ill{}} dans $\mathrm{GL}(r)$, $C$ dans $E^r$ \ldots{} ]

On trouve aussi que la donnée de $(P,\sigma)$ équivaut à celle d'une
\uncertain{structure d'}extension \ill{}

\page{2}
\note{ce feuillet est un début abandonné : il s'arrête au milieu d'une formule
et est repris, du même point de départ, au feuillet 3}
\[
  0 \longrightarrow \Omega \longrightarrow E \longrightarrow
  \underline{O}_S \longrightarrow 0
\]
en ce sens que le \uncertain{foncteur} qui à une telle extension fait
correspondre $(P,\sigma)$ avec $P = P(E)$ et $\sigma$ la section définie par le
quotient $\underline{O}_S$ de $E$, est une équivalence de catégories.

On trouve que les $\underline{O}_P/J^{\,n+1}$ (où $J$ est l'idéal de
$\sigma(S)$ dans $P$) s'identifient \struck{\ill{}} \marginal{aux
$\underline{O}_S$-algèbres augmentées à} $\mathrm{Sym}^n($ \ill{}

\page{3}
\note{le feuillet reprend, du même point de départ, le début abandonné du
feuillet 2}
\[
  0 \longrightarrow \Omega \longrightarrow E \xrightarrow{\;p\;}
  \underline{O}_S \longrightarrow 0
\]
$P = P(E)$, $\sigma$ section définie par $p$, $\widehat{P}$ complété formel de
$P$ le long de $\sigma$
\[
  \widehat{S} = \underline{O}_{\widehat{P}}
\]
$\Omega$ opère à g. sur $\widehat{S}$ --- car il opère sur l'extension (à g.)
donc sur $(P,\sigma)$ et $\widehat{P}$ (à droite)

$L = \underline{O}_P(1)$, $\widehat{L}$ module inv. sur $\widehat{S}$

$\Omega$ opère à g. aussi sur $\widehat{L}$, de façon compatible avec
opérations de $\Omega$ sur l'Algèbre de sections $\widehat{S}$.

Cas $E$ \struck{splittée} cohomologiquement triv. \ill{} --- on donne un
splitting $s : \underline{O}_S \to E$)
\[
  E^{\circ} \simeq \Omega \times \underline{O}_S
\]
Alors \struck{\ill{}}
\[
  \widehat{P}(S') \;\simeq\; \text{Quotients inversibles } L' \text{ de } E,
\]
qui sur $S$ réd. coïncident avec le quotient donné $\underline{O}_S$ --- donc
l'image de la section $e$ de $\underline{O}_S \subset E$ et une base de $L'$,
permettant d'identifier $L'$ à $\underline{O}_S$, de sorte que le quotient
\struck{soit} $\simeq \underline{O}_S$ et défini par
\[
  (x,\lambda) \longmapsto \varphi(x) + \lambda, \qquad
  \varphi : \Omega \longrightarrow \underline{O}_S
\]
bien mieux défini, de $\Omega$ dans \uncertain{le nilradical}
\marginal{$N_S$}

Donc
\[
  \widehat{P}^{\circ} = V(\Omega)^{\wedge} \quad
  \text{(complété formel le long de la section nulle)}
\]
\[
  = \mathrm{Spf}\;\widehat{\mathrm{Sym}}(\Omega)
  \;=\; \widehat{S}^{\circ}
\]
\note{à droite de cette ligne, une égalité que je lis
$\widehat{L}^{\circ} \simeq \widehat{S}^{\circ}$, l'exposant du second membre
restant douteux}

Il faut expliciter l'op. de $\Omega$ sur $\widehat{S}$. Soit
$a \in \Omega(S)$.

Donc $a$ opère sur l'ext. $E$ par
\[
  u_a(x,\lambda) = (x + \lambda a,\; \lambda)
\]
et si \struck{$\widetilde{\varphi} : E \to \underline{O}_S$}
$\varphi : \Omega \to N_S$, définissons le \ill{} $\widetilde{\varphi}$ de $E$,
$\widetilde{\varphi}(x,\lambda) = \varphi(x) + \lambda$, où avec le nouveau
quotient par
\[
  \widetilde{\varphi} \circ u_a : (x,\lambda) \longrightarrow
  \varphi(x + \lambda a) + \lambda
  = \varphi(x) + (1 + \varphi(a))\lambda
  = (1 + \varphi(a))\left( \frac{\varphi}{1 + \varphi(a)}(x) + \lambda \right),
\]
c'est aussi le quotient défini par
$(x,\lambda) \longmapsto \varphi^a(x) + \lambda$, donc

\page{4}
\struck{\ill{}} moyennant l'identification
$\widehat{P} \simeq V(\Omega)^{\wedge}$, les transports de l'op. $u_a$ donnent
sur $V(\Omega)^{\wedge}$
\[
  \varphi^a = u_a^{*}(\varphi) = \frac{\varphi}{1 + \varphi(a)}
\]
[ De \struck{\ill{}}, le morphisme induit $u_a^{*}(L') \simeq \underline{O}_S$
\ill{} sur les quotients inv. $L'$ de $\underline{O}_S$ (dans
$E = \Omega + \underline{O}_S$) est la multiplication par $1 + \varphi(a)$ ]

Donc, \struck{$u_a^{*}$} $\varphi$ est identifié à un homom. de
$\widehat{S} = \widehat{\mathrm{Sym}}(\Omega)$ dans $\underline{O}_S$
(coïncidant sur $\mathrm{Sym}^1$ avec celui qu'\ill{} d'ap.), on aura
\[
  \varphi^a(\omega) = \frac{\varphi(\omega)}{1 + \varphi(a)}
  = \varphi\Big( \sum_{n \geqslant 0} (-1)^n \omega a^n \Big),
\]
donc
\[
  \varphi^a(\omega) = \varphi(u_a(\omega))
\]
où
\[
  u_a(\omega) = \struck{\ill{}} \frac{\omega}{1+a}
  = \sum_{n \geqslant 0} (-1)^n \omega a^n
\]
\note{encadré à droite de cette ligne, la règle de degré :
$u_a(f) = f/(1+a)^n$ si $\deg f = n$}

\uncertain{Comme} est l'hom. can. épi
\[
  E \otimes_{\underline{O}_S} \widehat{S} \longrightarrow \widehat{L}
  \simeq \widehat{S}
\]
(quand $E \simeq \Omega + \underline{O}_S$, donc
$E \otimes_{\underline{O}_S} \widehat{S}
= \Omega \otimes_{\underline{O}_S} \widehat{S} + \widehat{S}$).
C'est-à-dire, restriction au facteur $\widehat{S}$ est
$\mathrm{id}_{\widehat{S}}$ ; sa restriction \ill{} s'identifie à un
$\underline{O}_S$-hom. $\Omega \to \widehat{S}$ qui implique $\Omega$ dans
$\widehat{S}^{+}$ --- d'où l'identité. Donc
\[
  \pi(\omega \otimes f,\, g) = \omega f + g,
  \qquad \omega \in \Omega(S),\; f,g \in \widehat{S}(S).
\]

Quel est l'endom. de $\widehat{L} \simeq \struck{\ill{}} \widehat{S}$
\uncertain{rendu} par l'autom. $u_a$ de la \uncertain{restriction} ? D'abord
trouvons l'autom. $u_a^E$ de $E \otimes \widehat{S}
= \Omega \otimes_{\underline{O}_S} \widehat{S} + \widehat{S}$ ; il est
$u_{a*}$ \ill{}
\note{la fin de la ligne, tracée sur la portée imprimée du feuillet, ne se
laisse pas lire}

\page{5}
en écrivant un élément de $E \otimes \widehat{S}$ (comme somme d'objets de ce
\ill{}) sous la forme
\[
  \omega \otimes f + e \otimes g,
  \qquad \omega \in \Omega(S),\;\; f,g \in \widehat{S}(S),
\]
$e$ étant la section \ill{}
\[
  u_{a*}^{\widehat{E}}(\omega \otimes f + e \otimes g)
  = u_a(\omega) \otimes u_{a*}(f) + u_a(e) \otimes u_{a*}(g)
\]
\[
  = \omega \otimes u_{a*}(f) + (e + a) \otimes u_{a*}(g)
\]
\[
  = \big( \omega \otimes u_{a*}(f) + a \otimes u_{a*}(g) \big)
  + e \otimes u_{a*}(g)
\]
On a donc
\[
  \pi\big( u_{a*}^{\widehat{E}}(\omega \otimes f + e \otimes g) \big)
  = \omega\, u_{a*}(f) + a\, u_{a*}(g) + u_{a*}(g)
\]
\[
  = \omega\, u_{a*}(f) + (a+1)\, u_{a*}(g)
\]
tandis que
\[
  \pi(\omega \otimes f + e \otimes g) = \omega f + g .
\]

\begin{tikzcd}
E \otimes \widehat{S} \arrow[r, "\pi"] \arrow[d, "u_{a*}^{\widehat{E}}"] &
  \widehat{S} \arrow[d, "u_{a*}^{\widehat{L}}"'] \\
E \otimes \widehat{S} \arrow[r, "\pi"] & \widehat{S}
\end{tikzcd}

Il faut avoir $\pi \circ u_{a*}^{\widehat{E}}
= u_{a*}^{\widehat{L}} \circ \pi$, et faire \ill{} --- \ill{}
\[
  u_{a*}^{\widehat{L}}(g) = (a+1)\, u_{a*}(g)
  \qquad \Big( = \frac{g}{(1+a)^{\,n-1}} \;\text{ si } \deg g = n \Big)
\]
\note{l'exposant du membre encadré se lit $n-1$, ce que la ligne précédente
demande}

Vérifions ceci \ill{}
\[
  u_{a*}^{\widehat{L}}\big( \pi(\omega \otimes f + e \otimes g) \big)
  = (a+1)\, u_{a*}(\omega f + g)
\]
\[
  = (a+1) \big( u_{a*}(\omega)\, u_{a*}(f) + u_{a*}(g) \big)
\]
\note{sous $u_{a*}(\omega)$, une accolade porte $\omega/(1+a)$}
\[
  = \omega\, u_{a*}(f) + (a+1)\, u_{a*}(g) \qquad \text{OK}
\]

§ On revient au cas général d'une extension quelc. de $\underline{O}_S$
par $\Omega$, défini par un \uncertain{torseur} $T$ sur $\Omega$ ; on trouve
donc
\[
  \widehat{S} = \struck{\ill{}}\;
  T \wedge^{(\Omega,\, u^{\widehat{L}}_{\Omega*})}
  \widehat{\mathrm{Sym}}(\Omega)
\]
\[
  \widehat{L} =
  T \wedge^{(\Omega,\, u^{\widehat{L}}_{\Omega*})}
  \widehat{\mathrm{Sym}}(\Omega)
\]
\note{les deux membres de droite sont écrits comme un produit contracté du
torseur $T$ par $\widehat{\mathrm{Sym}}(\Omega)$, le symbole de contraction
restant douteux}
\struck{\ill{}}
\note{les dernières lignes du feuillet sont biffées de traits diagonaux ; on
y lit $J^n/J^{n+1}$ et un quotient de $\mathrm{Sym}^n\Omega$ par
$\mathrm{Sym}^{n+1}$}

\page{6}
\note{brouillon au crayon, très pâle, traversé par la transparence du verso :
les formules se laissent lire, la prose qui les relie beaucoup moins}

\emph{Conséquence} : Soit $\widehat{J}$ \struck{\ill{}} l'idéal
d'aug\supplied{mentation} de $\widehat{S}$, qui \ill{}
\[
  \mathrm{gr}^1_{\widehat{J}}(\widehat{S}) \simeq \Omega
\]
d'où
\[
  \big( \mathrm{gr}_{\widehat{J}}(\widehat{S})
  \simeq \widehat{\mathrm{Sym}}(\Omega) \big)
\]
D'autre part, $J^n/J^{n+2}$ est une \ill{} ; \ill{}
$\mathrm{gr}^n \simeq \mathrm{Sym}^n(\Omega)$ \ill{}

\ill{} un calcul, on note que
$\mathrm{gr}^{\bullet}\,\widehat{J}^{(a,\omega)}/\widehat{J}^{\,\omega}$
\ill{}

Or $u_a$ \ill{} et donne pour \ill{}
\struck{$\sum \omega_i \cdots$}
\[
  f + g \;\longrightarrow\; \frac{f}{(1+a)^{n}} + \frac{g}{(1+a)^{n+1}}
\]
\note{sous $f$ et sous $g$ il porte leurs degrés, $n$ et $n+1$}
\[
  f \;\longmapsto\; naf + g + \text{termes de degré } \geqslant n + 1
\]
\note{le dernier terme de l'inégalité est douteux}

Soit donc
\[
  \Omega \;\longrightarrow\; \mathrm{Hom}\big( \mathrm{Sym}^n \Omega,\;
  \mathrm{Sym}^{n+1} \Omega \big)
\]
Ainsi \ill{} de
$\widehat{J}^{\,a,u} \big/ \widehat{J}^{\,a,u,u}$ \ill{}
\[
  v(a)(f) = af
\]
\struck{\ill{}}

(N.B. \ill{} $\Omega$ \ill{}.) Donc
\[
  T \times_{\Omega} \big[ \mathrm{Hom}\big( \mathrm{Sym}^n \Omega,\;
  \mathrm{Sym}^{n+m} \Omega \big) \big]
\]
\ill{} qui donne l'extension.

\page{7}
\note{feuillet au crayon, pâle ; les dernières lignes sont reprises à l'encre
bleue}

Soit $X$ schéma lisse sur $S$, \struck{\ill{}} et \ill{}

Une \emph{connexion projective} sur $X$ (rel. à $S$) est la donnée
\begin{itemize}
  \item[a)] d'un fibré de B.S. $P$ sur $X$ avec une section, \ill{} qui
  revient \ill{}, donc \ill{} d'une extension
  \[
    0 \longrightarrow \Omega \longrightarrow E \longrightarrow
    \underline{O}_X \longrightarrow 0
  \]
  de modules loc. libres sur $X$, \ill{}
  \item[b)] d'un $X$-isom. des voisinages \ill{} de la diagonale de
  $X \times_S X$ avec \ill{} dans $P$, \ill{} c'est-à-dire un isom. de
  $\underline{O}_X$-Algèbres augmentées
  \[
    P^n_{X/S} \;\simeq\; \underline{O}_{\widehat{P}}
    \big/ \widehat{J}^{\,n+1}
    \qquad \big( \simeq \underline{O}_X \otimes_{\Omega}
    \mathrm{Sym}^{\bullet}(\Omega) \big)
  \]
\end{itemize}
($\widehat{P}$ étant le complété formel de $P$ le long de \ill{} $(X)$, et
$\widehat{J}$ l'Idéal de \ill{}).

\ill{} la donnée de $P$ \ill{} revient \ill{}
\[
  (*) \qquad \Omega \;\simeq\; \Omega^1_{X/S}
\]
\note{$\Omega^1_{X/S}$ de $X$ sur $S$ : c'est l'identification que la suite
suppose}

Pour $n = 2$, la structure de $P^2_{X/S}$ \ill{} est déterminée par la
structure \ill{} d'extension de $\Omega^1_{X/S}$ par
$\mathrm{Sym}^2 \Omega^1_{X/S}$, \ill{} \struck{d'un cocycle} d'un torseur
sous \ill{} $\mathrm{Hom}_{\underline{O}}(\Omega^1_{X/S}, \ldots)$ \ill{}
doit être muni d'un iso \ill{}, lequel doit être \ill{} le torseur déduit du
$T$ sous $\Omega$ \ill{} par
\[
  \Omega \xrightarrow{\;-u_1\;}
  \mathrm{Hom}\big( \Omega,\ \mathrm{Sym}^2(\Omega) \big)
\]
et l'iso
\[
  \mathrm{Hom}\big( \Omega^1_{X/S},\ \mathrm{Sym}^2 \Omega^1_{X/S} \big)
\]
déduit de l'iso ci-dessus. Donc si l'on \ill{}

\page{8}
\ill{} de $X$ sur $S$ est $1$, \struck{\ill{}} \ill{} $u_1$, $-u_1$ \ill{}
$r+1$ \ill{} sur $S$ (\ill{} $u_1$, $-u_2$ \ill{} des \ill{} à $S$.),
\struck{le torseur} $T$ \ill{} déterminé par $X$, \ill{}
$\Omega \simeq \Omega^1_{X/S}$ et \ill{}

On dépend de \ill{} $=$ une \ill{} connexion (\ill{} n'a qu'une telle \ill{}
est can. déterminé par $X/S$, sans \ill{} donnée ni de con. ?)

Notons que $P^{\infty}_{X/S}$ a une \struck{plate} \ill{} canonique $\Delta$
\ill{}, donc \ill{} augmentation $P^{\infty}_{X/S} \to \underline{O}_X$ !),
donc si \ill{} d'une $\infty$-connexion projective \ill{} de $\widehat{P}$,
i.e. $\widehat{S}$ \ill{} \struck{\ill{}} correspondant.

On dit \ill{} qu'une $\infty$-conn. projective est \emph{intégrable} si elle
\ill{} pour une $\infty$-connexion \ill{} de $P = P(E)$ (un \ill{} par le
\ill{} a priori), laquelle sera dès lors \ill{} déterminée (tel qu'il \ill{}
« économique »).

Cette condition étant locale sur $X$, \ill{} l'exprimer en supposant que $X$
a \ill{} un \ill{} « coordonnées locales » $f_1, \ldots, f_n$ (\ill{} $n$ sur
$S$). Dès lors
\[
  P^{\infty}_{X/S} \simeq \underline{O}_S[[\delta f_1, \ldots]]
\]
et
\[
  \Delta : P^{\infty}_{X/S} \longrightarrow
  P^{\infty}_{X/S} \otimes_S P^{\infty}_{X/S}
  = P^{\infty}_{X/S}\big( P^{\infty}_{X/S} \big)
  \simeq P^{\infty}_{X/X}[[\ldots
\]
\[
  \Delta(\delta f_i) = \delta(f_i) \otimes 1 + 1 \otimes \delta(f_i)
  = \delta(f_i) + \delta f_i
\]
\[
  M \otimes P^n_{X/S} \longrightarrow P^n_{X/S} \otimes M
\]
\struck{\ill{}} $P^n_{X/S}$-linéaire \ill{}, i.e.
\[
  M \xrightarrow{\;\Delta\;} P^n_{X/S} \otimes M
\]
\ill{} $\underline{O}_S$-lin. par \ill{} ; id. \ill{} 2 termes
\uncertain{fantômes} \ill{}
\[
  \Delta(m) = m + \Sigma
\]

\note{la moitié inférieure du feuillet est partagée par un trait vertical ;
la colonne de droite porte : « On peut \ill{} l'extension $E$ \ill{}, donc
$\widehat{S} \simeq \widehat{\mathrm{Sym}}(\Omega)$, et comme
$\Omega \simeq \Omega^1_{X/S} \simeq \underline{O}^{\,n}$,
$\widehat{S} \simeq \underline{O}_S[[T_1, \ldots, T_n]]$ ;
[$\infty$-connexion projective \ill{}] »}

\page{9}
déterminé par les images des $T_i$ dans $P^{\infty}_{X/S}$, disons
\[
  F_i \in \Gamma\, P^{\infty}_{X/S}
  = \Gamma\, \underline{O}_S[[\delta f_1, \ldots, \delta f_n]],
  \qquad 1 \leqslant i \leqslant n
\]
\[
  F_i = \delta f_i + G_i, \qquad G_i \text{ d'ordre } \geqslant 2
\]
ou inversement par les images des $\delta f_i$, soit
\[
  \Phi_i \in \Gamma\, \widehat{\mathrm{Sym}}\;
  \underline{O}_S[[T_1, \ldots, T_n]],
  \qquad 1 \leqslant i \leqslant n
\]
\[
  \Phi_i = T_i + \Psi_i \qquad (\Psi_i \text{ d'ordre } \geqslant 2)
\]
On a donc \struck{la structure} \ill{} sur $\underline{O}_S[[T_1, \ldots$
\[
  \underline{O}_S[[T_1, \struck{\delta f}]] = \widehat{S}
  \xrightarrow{\;\Delta'\;} P^{\infty}_{X/S} \widehat{\otimes} \widehat{S}
  = P^{\infty}_{X/S}(\widehat{S}) \simeq P^{\infty}_{X/S}[[T \ldots
\]
déduit par transport de structure de celle de $P^{\infty}_{X/S}$. Il faut
donner les i\ill{} des $T_i$ dans $P^{\infty}_{X/S}[[T_1, \ldots, T_n]]$,
\ill{} l'image de
\[
  \Delta(\struck{\ill{}}\; \delta f_i + G_i) \in \Gamma\, P^{\infty}_{X/S}
\]
de $\Gamma\, P^{\infty}_{X/S} \widehat{\otimes} \widehat{S}$, déduit de
$\longrightarrow$ obtenu par \ill{}
\[
  c_{\infty} : P^{\infty}_{X/S} \longrightarrow \widehat{S}.
\]
Posant
\note{le feuillet s'arrête au premier tiers, sur ce mot ; la suite manque}

\page{10}
\begin{tikzcd}
X \arrow[d, "f"] \\ S
\end{tikzcd}

$X$ propre lisse dim. rel. $1$, fibres connexes
\[
  \Omega^1_{X/S} = \omega, \qquad f_*(\omega) = \Omega, \qquad
  R^1 f_*\big( \Omega^{\bullet}_{X/k} \big) = DR
\]
\[
  R^1 f_*(\underline{O}_X) \simeq \check{\Omega} \simeq \tau
\]
\note{le troisième terme est douteux ; je le lis $\tau$}
\[
  0 \longrightarrow \Omega \longrightarrow DR \longrightarrow
  \check{\Omega} \longrightarrow 0
\]
\marginal{autodual}
\struck{\ill{}}
\[
  DR \times DR \xrightarrow{\;q\;} \underline{O}_S \qquad \text{alternée}
\]
\[
  b = R^2 f_*\big( \Omega^{\bullet}_{X/S} \big)
  \simeq R^1 f_*\big( \Omega^1_{X/S} \big)
\]
\[
  0 \longrightarrow \check{\Omega} \otimes \check{\Omega} \longrightarrow
  \Sigma \longrightarrow \underline{O}_S \longrightarrow 0
\]
\ill{} le groupe de $\Omega^2 \to \Omega$, \struck{$\Omega \otimes
\check{\Omega}$}
\[
  \Lambda^2 DR \xrightarrow{\;q\;} \underline{O}_S
\]

\begin{tikzcd}
0 \arrow[r] & \Omega \arrow[r] \arrow[d, "\mathrm{id}"] &
  DR \arrow[r] \arrow[d, "2"] &
  \check{\Omega} \arrow[r] \arrow[d, "\mathrm{id}"] & 0 \\
0 \arrow[r] & \Omega \arrow[r] & DR \arrow[r] &
  \check{\Omega} \arrow[r] & 0
\end{tikzcd}
\note{l'étiquette de la flèche médiane se lit $2$, sans certitude}

\ill{} on considère \ill{} l'accouplement \ill{} $\Omega \times DR$ et
$DR \times \Omega$ --- l'aboutissement \ill{} dans \struck{\ill{}}
$\mathrm{Hom}(\check{\Omega} \otimes \check{\Omega}, \underline{O}_S)
= \Omega \otimes \ldots$ \ill{} $= \ell$ puisqu'il s'agit de formes alternées
\ill{} dans $\Lambda^2 \Omega$, $\Omega \otimes \check{\Omega}$,
$\Lambda^2 \check{\Omega}$
\[
  0 \longrightarrow \mathrm{Fil}^2 \longrightarrow \mathrm{Fil}^1
  \longrightarrow \Lambda^2 DR
\]
\note{une flèche portant « inv » descend de $\mathrm{Fil}^1$ vers
$\underline{O}_S$}

Pour un splitting $q$ de $DR$ \struck{\ill{}} $q : \Omega \to DR$ \ill{} on
aura comme matrice de toute p.c. la décomposition
$DR \simeq \Omega \oplus \Omega$
\[
  \begin{pmatrix} 0 & \pm\,\mathrm{id} \\ \mathrm{id} & \varphi_0 \end{pmatrix}
\]
\ill{} pour \ill{} déduite sur $\check{\Omega}$. P\ill{} un autre choix
\[
  q'(t) = q(t) + u(t) \qquad \big( u : \check{\Omega} \to \Omega \big)
\]
\ill{} on aura
\[
  \varphi'_0(t,t') = \varphi\big( q(t) + u(t),\; q(t') + u(t') \big)
\]
\[
  = \varphi_0(t,t') + \big( t,\, u(t') \big) - \big( t',\, u(t) \big)
\]
\note{le signe entre les deux derniers termes est masqué par un pâté d'encre ;
je le lis $-$}

\page{11}
\note{feuillet à l'encre bleue, d'une écriture rapide : les formules portent,
la prose qui les relie est en grande partie perdue}

Si \ill{} l'inverse \ill{} où le $u$ \ill{} $\varphi'_0 = 0$, on voit que
l'on a une condition \ill{} et \ill{} torsion \ill{} $\Lambda^2 \Omega$
\ill{} torseur sous \ill{} (formes symétriques sous $\Omega$ \ill{})
\[
  0 \longrightarrow \Lambda^2 \Omega \longrightarrow \Sigma \longrightarrow
  \underline{O}_S \longrightarrow 0
\]
(Il \ill{} induit par un torseur \ill{} hyperbolique \ill{} les torseurs sous
$\pi_1 \subset \mathrm{Aut}$ \ill{} les schémas \ill{} pour les \ill{}
problème ?)

On a comme \ill{} (si $\Sigma$ \ill{} $T$)
\[
  DR \longrightarrow DR \otimes \Omega^1_{S/T}
\]
\ill{} compatible avec la filtration \ill{} ; les défauts de compatibilité
sont précisés \ill{}
\[
  \Omega \longrightarrow \check{\Omega} \otimes \Omega^1_{S/T}
\]
\ill{}
\[
  \Omega^1_{S/T} \longrightarrow \check{\Omega} \otimes \check{\Omega}
\]
\ill{} est \struck{\ill{}}, induit par le foncteur \ill{} pour problème
\ill{} et sur le contexte \ill{} par le \struck{\ill{}} un iso ; ici donné,
dont \ill{} par les conditions \ill{} $S/T$ \ill{} hyperbolique), \ill{} qui
\ill{}
\[
  \Omega^1_{S/T} \longrightarrow \Lambda^2 \check{\Omega}
\]
(\ill{} pour les \ill{} avec vol. \ill{} permutation avec \ill{})

\ill{} Sur $X$ (de genre $g \geqslant 2$) \ill{} fibré \ill{} d'un
sous-fibré de corang $1$ (de \ill{})
\[
  H^0(X, \Omega_X)
\]
(\struck{\ill{}} \ill{} formes \ill{}) s'annulant \ill{}
\[
  0 \longrightarrow \Omega^1_X \xrightarrow{\;\omega\;} \Omega_X
  \xrightarrow{\;\det\;} \Omega^1_X \longrightarrow (DR)_X
\]
de $DR$,
\note{la lecture de cette dernière suite est incertaine dans le détail des
flèches}

\page{12}
\struck{splittings} \marginal{$\varphi$} splittings de $DR$ $\longrightarrow$
formes \ill{} symétriques \ill{} sur $DR$
\[
  0 \longrightarrow \Gamma^2 \Omega \longrightarrow \Omega \otimes \Omega
  \longrightarrow \Lambda^2 \Omega \longrightarrow 0
\]
\note{une flèche descendante part de $\Gamma^2\Omega$, étiquetée « Epi (car
\ill{}) (car. $\neq 2$) »}
\[
  f_*(\omega^2) \;\simeq\; \Omega^1_{S/T}
\]
\[
  \frac{g(g+1)}{2} - (3g-3) \;=\; \frac{1}{2}\big( g^2 - 5g + 6 \big)
  \;=\; \frac{1}{2}(g-2)(g-3)
\]
\note{sous le membre médian, une accolade porte $(g-2)(g-3)$}
\[
  \longrightarrow \underline{O}_X \longrightarrow P^1 \longrightarrow \omega
  \longrightarrow 0
\]
\[
  E \otimes_{\underline{O}} P \longrightarrow E
\]
\[
  2 - 2g
\]
\[
  P^{\infty}_{X/S}, \qquad \Omega\quad \Omega^2\quad \Omega^3
\]
\[
  \omega^2 \longrightarrow P^{3+} \longrightarrow \omega \longrightarrow 0
\]
\[
  \longrightarrow f_*(\omega^2) \longrightarrow f_*\big( P^{3+}_{X\varphi}
  \big) \longrightarrow f_*(\omega) \longrightarrow 0
\]
\note{sous les deux membres extrêmes il écrit $\Omega^1_{S/T}$ et $\Omega$ ;
à droite, $J \subset A \otimes A \to A$, puis $J/J^2$ et
$x \otimes y - y \otimes x$}
\[
  P^{\infty}\big( \widehat{\mathrm{Sym}}(\Omega) \big), \qquad
  P^{\infty}(\omega^n), \qquad \omega \otimes P^{\infty},
  \qquad P^{\infty} \otimes \omega
\]
\[
  M \otimes P^{\infty}, \qquad R^1 f_*\big( M \otimes \omega^n \big) = 0
  \quad \text{si } n \gg 1
\]
\[
  S^i = f_*\, P^{\infty}(\omega^n) = f_*\big( \omega^n \otimes P \big)
\]
\[
  \mathrm{gr}\, S^{\bullet} \;\simeq\; \sum_{n \geqslant \ill{}}
  f_*\big( \omega \otimes P \big).
\]
\note{la dernière ligne porte un indice de sommation que je ne lis pas ; à
gauche, un petit diagramme $\Sigma^{\bullet} \to S^{\bullet}$ avec une flèche
montante depuis $R^{\bullet}$, et $\Sigma \otimes P^{\bullet}$}

\page{14}
\note{le feuillet 13 ne porte que le numéro des archivistes}
\note{à partir d'ici il note en exposant un signe que je ne sais pas nommer,
et que je rends par $\mathfrak{z}$ ; c'est le même qui donne
« $\mathfrak{z}$-structure » dans la prose, et il y désigne la structure
projective}

Fixons \ill{} un genre $g$ ($> 2$ ou \ill{}) et un foncteur rigidifiant
(p. ex. de Torelli, d'i\ill{}els \uncertain{convenables}) pour les courbes
algébriques \ill{} relatives sur une base donnée (\ill{} ou travaillant sur
la \struck{\ill{}} \marginal{multiplicités} \ill{} modulaire des courbes de
genre $g$ --- lisse, propre, groupe \ill{}). La catégorie fibrée qui s'\ill{}
ou la \ill{} \ill{}, soit notée $M_g$.

Considérons
\[
  M_g^{\mathfrak{z}}(S) = \text{courbes algébriques}
\]
\marginal{$S$ schéma de base de car. $0$}
propres lisses $X$ de genre $g$ sur $S$, munies d'un \ill{} fibré par des
\ill{} projectives $P$ sur $X$, avec une $\mathfrak{z}$-structure (connexion
lisse) i\ill{} sur $S$, telle que $\forall\, s \in S$ \ill{} de $S$,
$P_s$ sur $X_s$ ait un \uncertain{point} $\mathfrak{z}$-ordinaire
\ill{} le groupe structural de $P$ \ill{} $\Pi_0 \subset \mathrm{PGL}(1)$,
\ill{} $\subset \mathrm{PGL}(1)$.

(i.e. $P$ n'est \uncertain{intégrable} : \ill{} « dégénéré » \ill{})

\emph{Rq.} Soit $X$ \struck{courbe} \marginal{schéma} algébrique connexe sur
$k$ alg. clos, \struck{\ill{}} $P$ un $\mathfrak{z}$-fibré ou élément
projectif sur $X$, conditions équivalentes :
\begin{itemize}
  \item[a)] $\mathrm{Aut}^{\mathfrak{z}}(P) = 1$
  \item[b)] $\mathrm{Aut}^{\mathfrak{z}}_{X/k}(P)$ (qui est un groupe alg.
  $\subset \mathrm{PGL}(1)_k$) est de dim. $0$
  \item[c)] $P$ n'est pas de l'un des deux types
\end{itemize}
\begin{itemize}
  \item[$\alpha$)] associé à un \ill{}-$\mathfrak{z}$-fibré, i.e. à une
  $\mathfrak{z}$-\ill{} isomorphe \struck{\ill{}}, un point, la clôture
  projective de $V(\Omega)$, i.e. $P(\Omega + \underline{O})$
  \item[$\beta$)] associé à un $\mathbf{G}_a$-$\mathfrak{z}$-fibré, i.e. à
  une $\mathfrak{z}$-extension
  \[
    0 \longrightarrow \underline{O}_X \longrightarrow E \longrightarrow
    \underline{O}_X \longrightarrow 0
  \]
  (\ill{} $\mathfrak{z}$-structures triviales sur $E$) ou plutôt
  $P(E) \simeq$ clôture projective du torseur \ill{}
\end{itemize}
\marginal{l'intersection de ces deux types de $\mathfrak{z}$-structures (i.e.
$\mathfrak{z}$ triviales incluses)}

\page{15}
\ill{} Si $X$ courbe algébrique \marginal{propre} lisse \ill{} de genre $g$
\ill{} \struck{\ill{}}
\begin{itemize}
  \item[d)] $\dim H^1_{DR}(X, \mathfrak{g}) = 6g - 6$ \struck{\ill{}}
  (\ill{} jamais vérifiée si $g \geqslant 0$ --- et bien sûr !)
  \item[d')] $H^0_{DR}(X, \mathfrak{g})$ $\big( = H^0(X, \mathfrak{g})\big)
  = 0$
  \item[d'')] $H^2_{DR}(X, \mathfrak{g})$ $\big( = H^2(X, \mathfrak{g})\big)
  = 0$
\end{itemize}

[N.B. On a $\mathfrak{g} \simeq \check{\mathfrak{g}}$ (Killing), donc
\[
  (*) \qquad H^0_{DR}(X, \mathfrak{g}) \simeq
  H^2_{DR}(X, \mathfrak{g})^{\vee}
\]
(d'où équivalence de d) et d''\,)) et formule d'E.P.
\[
  (**) \qquad \chi_{DR}(X, \mathfrak{g})
  \overset{\text{déf}}{=} \sum (-1)^i \dim H^i_{DR}(X, \mathfrak{g})
  \;=\; 3\, \chi_{DR}(X, \underline{O}_X)
\]
\note{sous $\chi_{DR}(X,\underline{O}_X)$ il porte sa valeur, $2-2g$}
\[
  = \struck{6g} \; 6 - 6g
\]
donc, combinant avec $(*)$,
\[
  (\!*\!*\!*) \qquad \dim H^1_{DR}(X, \mathfrak{g})
  = 2 \dim H^0_{DR}(X, \mathfrak{g}) + (6g - 6)
\]
d'où l'équivalence de d'\,) avec d). Enfin, on sait
\[
  H^0_{DR}(X, \mathfrak{g}) \simeq \mathrm{Lie}\;
  \mathrm{Aut}^{\mathfrak{z}}_{X/k}(P)
\]
d'où l'équivalence de d'\,) avec b), et la formule \ill{}
\[
  \dim H^1_{DR}(X, \mathfrak{g}) = 2 \dim
  \mathrm{Aut}^{\mathfrak{z}}_{X/k}(P) + 6g - 6
\]
\[
  = 2 \Big( \dim \mathrm{Aut}^{\mathfrak{z}}_{X/k}(P) + 3g - 3 \Big)
\]
\note{une accolade sous $\dim H^1_{DR}(X,\mathfrak{g})$ porte : « espace
tangent à l'espace modulaire $M_g^{\mathfrak{z}}$ au pt correspondant à
$(X, P, \mathfrak{z})$ »}

N.B. La condition d'\,) implique que $M_g^{\mathfrak{z}}$ « n'a pas
d'automorphismes infinitésimaux », c) qu'il n'a pas d'automorphismes finis »,
donc on \ill{} pas $M_g^{\mathfrak{z}}$ est « rigide » --- tandis que comme
$M_g$ a \ill{} dans un l'espace représentable ! Le semble. d''\,) implique
que $M_g^{\mathfrak{z}}$ est \marginal{formellement} lisse sur $M_g$, et d)
que sa dim relative sur $M_g$ est $6g - 6$.
\note{la dimension relative se lit bien $6g-6$}

\marginal{Sans doute peut-on prouver la représentabilité \ill{} relative
(dans un cas \ill{} \uncertain{org. d'Abel} \ill{}}
\note{cette remarque est portée en haut de la moitié droite du feuillet et
court jusqu'au bord, où elle est coupée}

\page{16}
$X$ courbe analytique connexe compacte \ill{},
$M_g^{\mathfrak{z}}(X)$ espace modulaire pour les \ill{} projectifs
analytiques : connexions plates « sans automorphismes ». C'est l'espace
analytique associé à une solution (de \ill{}) sur $\mathbf{C}$,
\struck{mais} lisse de dim. $6g-6$ \ill{} --- le \uncertain{regard} \ill{}
de $M$ \ill{} analytique. Il est donc \ill{} (\ill{} pas choisi qu'un bon
$x \in X$) \ill{} des représentations
\[
  \pi \xrightarrow{\;\rho\;} \mathrm{PGL}(1, \mathbf{C}) = H,
  \qquad \pi = \pi_1(X, x)
\]
\struck{telles que} \ill{} telles que $\mathrm{Cent}(\varphi) = \{e\}$,
\struck{ou} qu\ill{} $H$ \ill{} intérieurs
\[
  M_g^{\mathfrak{z}}(X) \;\simeq\;
  \mathrm{Hom}_{\mathrm{gr.\,ana}}(\pi, H) \,\big/\, H
\]
\ill{} le \,!\, désigne le \ill{} des
$\varphi \in \mathrm{Hom}_{\mathrm{gr}}(\pi, H)$ dont le st\ill{} dans $H$ se
réduit à $1$.

[ \ill{} de faits et \ill{} une \ill{} algébrique de \ill{} $g \in \pi$, le
\ill{} \ill{} algébrique), et \ill{} algébrique. Mais l'iso \ill{}
\uncertain{sans doute} \ill{} algébrique ! ] L'interpr\ill{} \ill{} de \ill{}
classifie les couples $(P, \mathfrak{z}, \varphi)$, $(P, \ldots$

\page{17}
\note{le feuillet est écrit sur deux colonnes ; la colonne de droite est
lourdement raturée et sa moitié inférieure, surchargée de renvois, ne se
laisse pas suivre}

$\mathfrak{z}$-fibré en droites projectives, et
$\varphi : P(\mathfrak{z}_1) \to P_2$ un iso. L'opération de $H$ dessus est
l'évidente.

Or le groupe $\mathrm{Aut}(\pi)$ opère \marginal{à droite} sur
$\mathrm{Hom}(\pi, H)$ ; si $u \in \mathrm{Aut}(\pi)$,
$\rho \in \mathrm{Hom}(\pi, H)$, \ill{}
\[
  \rho \cdot u = \rho \circ u ;
\]
si $u = \mathrm{int}(g)$ ($g \in \pi$), \ill{}
\[
  \rho \circ u = \mathrm{int}(\rho(g)) \circ u,
\]
donc \ill{} un groupe au quotient : $\mathrm{Hom}^{\,!}(\pi, H)/H$,
c'est-à-dire $\gamma : \struck{\ill{}}\ \mathrm{Out}(\pi)
= \mathrm{Aut}(\pi)/\mathrm{Int}(\pi)$, qui opère sur lui. Par transport de
structure, ce groupe opère sur $M_g^{\mathfrak{z}}(X)$, mais alors dont cette
opération soit \uncertain{algébrique}.

Si $Y$ est une courbe analytique complète \struck{de} genre $g$, alors la
donnée d'un iso. est
\[
  \pi_1(Y) \;\simeq\; \pi = \pi_1(X, x)
\]
définit, si la représentation composée
\[
  \rho_Y : \pi_1(Y) \longrightarrow H
\]
(définie \ill{} à \ill{} int. dans $H$), un élément de
$M_g^{\mathfrak{z}}(X)$. On trouve ainsi, si $M_g(\pi) \struck{= M_g(X)}$
\ill{} l'espace de Teichmüller (à rigidification définie par $\pi$), \ill{}
trouve un \uncertain{morphisme} analytique \marginal{injectif}
\[
  M_g(X) \hookrightarrow M_g^{\mathfrak{z}}(X)
\]
qui commute aux \ill{} $\mathfrak{z}$ \ill{} des groupes de

\note{colonne de droite}
Teichmüller $\Gamma$. On suppose que l'image \struck{\ill{}}
\marginal{\ill{} compact} par \ill{} constructible \ill{} algébrique. Elle
est formée des $\mathrm{Hom}$ injectifs $\pi \subset \ell$, $H$ pour
lesquels il existe un domaine hom\ill{}gique
$D \subset \mathbf{P}^1(\mathbf{C})$ \ill{} par $\rho$ (il sera unique) et
que \struck{$D/\pi$ \ill{}} \struck{\ill{} opère librement} et \ill{} un
quotient compact.

On définit aussi $M_{\mathrm{ann}}(X)$ qui classifie les $P(\ldots)$
\struck{\ill{}} \ill{} sur $X$, munis d'une section nulle \ill{}
dictionnaire. \struck{Il est} (c'est un \ill{} sur le \ill{} de deux fois
\[
  H^0\big( X, \omega^{\otimes 2}_{X/\mathbf{C}} \big), \ldots
\]

\[
  M_{\mathrm{ann}}(X) \longrightarrow M_g^{\mathfrak{z}}(X)
\]
\ill{} est un morphisme \uncertain{net}. (Est-il injectif ?) (N.B. C'est en
fait un morphisme de \ill{})

\page{18}
\ill{} classe de torsion dans \uncertain{$\mathrm{Pic}$}, $L \otimes
\omega^{-1}$ anti-ample. Donc tout aut. dét\ill{} de
$P\big( P^1_{X/S}(-1) \big)$ laissant fixe le \ill{}
\note{le tiers supérieur du feuillet est ensuite annulé par des hachures ;
on y distingue encore $M_g^{\mathfrak{z}}(X)$ et $M_g(X)$}

\ill{} $L \simeq \underline{O}_X$, donc \ill{} un fait \ill{} autre $\omega$
de $E$, avec $u(\omega) \subset \omega$, et \ill{} \struck{Plus} le quotient
$E/\omega \simeq \underline{O}_X$. Si \ill{} $\varrho_\alpha$ \ill{}, est
induit \ill{} $\lambda \in H^0(X, \underline{O}_X^{*}) = \mathbf{C}^{*}$ sur
$\omega$, \struck{l'extension} $c \in H^1(X, \underline{O})$ \ill{}
$\mathbf{C}$ dans \ill{} l'ex. $E$ de $\underline{O}_X$ par $\omega$, \ill{}
donc avoir
\[
  \lambda c = c \quad\text{i.e.}\quad (\lambda - 1)\, c = 0,
\]
et comme $c \neq 0$ \ill{} doit avoir $\lambda = 1$. Donc \ill{} un
automorphisme d'extension, donc défini \ill{} un homom.
$\underline{O}_X \to \omega$, i.e. par une forme diff.
$\omega \in H^0(X, \omega)$. Donc
\[
  \mathrm{Aut}_X(P_X) \;\simeq\; H^0(X, \omega_X)
\]
\note{au-dessus du membre de gauche il écrit $\mathrm{Aut}_X(P_X, \sigma)$,
et sous l'isomorphisme la correspondance $u_\omega \leftrightarrow \omega$}
(iso. de groupes). La question est alors la suivante : peut-on avoir
$u_\omega(c) = c^{\ill{}}$ si $c$ \struck{\ill{}} \marginal{est la} classe
sur $P_X$ associée \ill{} \struck{deux} classe projective $C$,
\struck{\ill{}} ?

\page{19}
Choisissons une uniformisante \struck{\ill{}} $t$, d'\ill{} \ill{} $\omega$,
d'où \uncertain{splittages} de $\pi$ et de $\mathcal{U}$
\[
  \mathcal{U} \simeq \omega + \underline{O}_X + \omega^{-1}
\]
et connexion projective $c_{D_t}$ donnée par
\[
  c_{D_t} : \mathcal{U} \longrightarrow \mathcal{U} \otimes \omega
  \simeq \omega^2 + \omega + \underline{O}_X
\]
\[
  \begin{cases}
    c_{D_t}(dt) = -2\, dt \\
    c_{D_t}(1) = 1 \\
    c_{D_t}\big( (dt)^{-1} \big) = 0
  \end{cases}
\]
donc
\[
  c_{D_t}\big( \lambda\, dt + \mu \cdot 1 + \nu\, (dt)^{-1} \big)
  = \lambda'\, dt^2 + (\mu' - 2\lambda)\, dt + (\nu' \ldots
\]
Soit $\omega \in \Gamma(\omega)$, d'\ill{}
\[
  u_\omega = \exp \Theta_\omega : \mathcal{U} \longrightarrow \mathcal{U}
\]
et c\ill{} que
\[
  \exp \Theta_{-\omega}(1) = 1 + [-\omega,\, 1] = 1 + \omega
\]
On veut calculer :
\[
  u_\omega(c_{D_t}) = \struck{\ill{}}\; u \circ c_{D_t} \circ u^{-1}
  = \big( \exp \Theta_{-\omega} \big) \ldots
\]
Et pour ceci il suffit de calculer
\[
  u_\omega(c_{D_t}) - c_{D_t} \;\equiv\; \varphi
\]
\ill{} comme \ill{} $\varphi(dt)$ \ill{}
\[
  \varphi(dt) \;\struck{=}\; c_{D_t}(dt)
  = \big( \mathrm{id} + \Theta_{-\omega} \big)
  \big( \struck{D_t}(dt) \big) = -2\, dt + \ldots
\]
\[
  \varphi(1) = -2\,\omega\, dt
\]
\[
  \varphi(1) + c_{D_t}(1) = \Big( \mathrm{id} + \Theta_{-\omega}
  + \tfrac{1}{2} \Theta_{-\omega}^{\,2} \Big)\,
  c_{D_t}\big( 1 + [\omega,\, 1] \big)
\]
\note{sous $c_{D_t}(1)$ il écrit $1$ ; sous $1 + [\omega,1]$, $1 - \omega$}
\[
  = 1 + [-\omega,\, 1] + \big[ -\omega,\, [-\omega,\, 1] \big] + \ldots
\]
\note{sous $[-\omega,1]$ il écrit $-2\omega$, et sous le crochet intérieur
$-2\omega$ également}
\[
  = 1 + (-2\omega) + 2\omega^2 - \lambda'\, dt^2 + 2\lambda\, dt + \ldots
\]
\[
  \varphi(1) = \big( -\lambda' + \struck{\ill{}} \big)\, dt^2
  \;\struck{\ill{}}
\]
\note{la dernière ligne du feuillet est entièrement biffée}

\page{20}
On trouve
\[
  u_\omega(c_{D_t}) = c_{D_t} + d^1_{\omega/S}(\omega)
\]
et plus généralement
\[
  \boxed{\;u_\omega(c) = c + d^1_{\omega/S}(\omega)\;}
\]
si $c$ connexion projective.

Cela montre :
\begin{itemize}
  \item[a)] $u_\omega(c)$ n'est une connexion projective \ill{} ssi
  $d\omega = 0$, i.e. (si $2$ inv. dans $S$) $\omega = 0$ : donc
  \struck{pas} \marginal{et rien pas} deux connexions projectives
  différentes ne peuvent donner lieu \ill{} du $\mathfrak{z}$-fibré. Donc
  \ill{}
  \item[b)] On n'a $u_\omega(c) = c$ que si $d^1_{\omega/S}(\omega) = 0$,
  donc a fortiori $d\omega = 0$ \ill{} (si $2$ inv. dans $S$) \ill{} il
  \ill{} de $\mathfrak{z}$-fibrés \ill{}
\end{itemize}

\note{suit un bloc entièrement biffé de traits obliques, où se lisent
$M_g(X)$, $M_g^{\mathfrak{z}}(X)$ et $\mathrm{conn}(X)$ reliés par des
flèches}

Non, car s'il y en avait, le $\mathfrak{z}$-\ill{} serait
\uncertain{énoncé} : un $\mathfrak{z}$-fibré \uncertain{affine} de \ill{}
serait isomorphe : un $P(F)$, où $F$ \ill{} $\mathfrak{z}$-extension
\[
  \ill{} \longrightarrow M \longrightarrow F \longrightarrow
  \underline{O}_X \longrightarrow 0
\]
\ill{} $M$ un $\mathfrak{z}$-fibré \ill{} inv. \ill{} ; on aurait donc
\[
  E \simeq F \otimes L,
\]
et \struck{\ill{}}

\end{document}
